\begin{center}
ABSTRACT
\end{center}

\bigskip

\begin{center}
SIRENS OF THE SWARM:\\
REVEALING BINARY BLACK HOLE MERGERS AND SUPERNOVA IN ACTIVE GALACTIC NUCLEI WITH OBSERVATIONAL PREDICTIONS FOR THE AGE OF MULTIMESSENGER ASTRONOMY \\
\bigskip
BY\\
\bigskip
HARRISON EDWARD COOK, B.A.
\end{center}

\bigskip

\begin{center}
Doctor of Philosophy\\
New Mexico State University\\
Las Cruces, New Mexico, 2022\\
Dr. Wladimir Lyra, Advisor
\end{center}

\bigskip
\bigskip

% ABSTRACT (350 WORDS)
%The atmosphere of Jupiter hosts a complex array of dynamical phenomena, which are tied to the planet's chemical inventory.  Observations from the \textit{Juno} spacecraft have provided various observational constraints on the various dynamical processes that occur on the planet such as the myriad of storms dominating the polar regions, possibly deep seated jets in the lower latitudes, and a large catalogue of atmospheric waves that are ubiquitous throughout. The observed dynamics are a function of deep atmospheric behavior, particularly where thermochemical activity of the atmosphere becomes relevant.  Thus, the current state of Jupiter is a function of its chemical constituents, the active role they play in shaping the atmospheric circulation, and the history of the giant's formation.
% Observations from ground-based telescopes to spacecraft alike, show a vast gaseous landscape of immense complexity.
\par
%Observations of trace disequilibrium chemical species are thought to be the result of atmospheric vertical mixing.  1D models, which ignore most of the dynamics in favor of chemical complexity, are able to provide an estimate of deep atmospheric abundances of such species, constraining the overall chemical inventory of the planet, particularly oxygen.  However, there is no consistent model that ties together the behavior of deep atmospheric chemistry to the observed atmospheric dynamics.  The abundance of water in the deep atmosphere has gross chemical consequences and shapes the nature of moist convective events, as seen in the tropospheric regions.  The purpose of this research is to connect the disequilibrium chemical behavior of the Jovian troposphere and deeper atmosphere to its moist convective and diffusive processes.
% Chromospheric signatures of nonthermal flare energy deposition are varied, and include small-scale brightenings around the energy injection sites, strong velocity flows that are frequently observed to be bidirectional, dramatic increases in the widths of spectral lines, and significant density enhancements. The purpose of this research is to connect the profile of nonthermal energy injection to the deposition and response throughout the chromosphere and transition region.
\par
%For the purposes of studying chemical behavior near the water cloud decks, the SNAP hydrodynamics model is used.  This model is non-hydrostatic and includes phase change and condensation effects.  Using SNAP, we model the behavior of carbon monoxide (CO) in Jupiter's atmosphere to constrain the effect of molecular weight changes at the cloud levels.  This formulation allows for an exploration of how CO abundance is effected by the Jovian oxygen inventory, as the oxygen content was found to directly impact the vertical mixing efficiencies due to cloud formation.  For the purposes of studying the diffusive properties of the deeper atmosphere and the corresponding vertical mixing coefficient at depth, we develop an analytical formulation of the eddy viscosity and test its applicability using VULCAN, a 1D chemical-diffusion code.
\par
%This work aims to synthesize the overall picture of the Jovian atmosphere into a cohesive story.  The first part of the work addresses the atmosphere near the cloud decks, where moist convective processes are numerous and lead to varying CO distributions.  Whereas, in the deeper atmosphere, the presence of a stably stratified layer may change the nature of the vertical diffusion.
% For the purposes of studying flare-driven chromospheric evaporation, a study is made of an X1.6 flare, SOL2014-10-22. This entirety of the nonthermal injection event was covered by a variety of instruments. From this, we derived the energy contained in nonthermal electrons, and link the time-dependent profile of energy injection to the response in Doppler velocity, nonthermal velocity, and density, as a function of time, temperature, wavelength and space. The fast-temporal variations in nonthermal energy injection and deposition is studied via ground-based observations of the low chromosphere, which are connected to X-ray flux as a proxy for nonthermal energy. These observations provide invaluable observations of fundamental time variations of the nonthermal event, and connect these to the smallest-scale X-ray linked structures. Taken together, this work provides valuable characterizations and constraints for the refinement, guidance, and initialization of the hydrodynamic flare simulation codes that will be valuable tool for advancing our understanding of the flare phenomenon throughout the peak of the next solar cycle.
